\begin{frame}{Коннектом ограничивает, но не определяет нейродинамику}
\begin{itemize}
\item Основной источник --- взрослая Drosophila melanogaster. Личиночные цепи будут анализироваться отдельно.
\item Коннектом и функциональные наблюдения будут согласованы по версии, стадии развития и исследуемой цепи.
\item Генератор будет сохранять воздействие, нейронные состояния, начальные условия, единицы измерения и параметры случайной инициализации; будут проверены воспроизводимость и чувствительность к топологии.
\end{itemize}
\SlideGap
Граф связности не задаёт временные свойства активности однозначно: нужны уравнения состояния, параметры взаимодействия и независимые наблюдения. Результаты моделей сенсомоторных и зрительных цепей сами по себе не обосновывают перенос к ECAP человека.
\SourceNote{Beiran, Litwin-Kumar, Nature Neuroscience, 2025. DOI: \href{https://doi.org/10.1038/s41593-025-02080-4}{10.1038/s41593-025-02080-4}. Shiu et al., Nature, 2024. DOI: \href{https://doi.org/10.1038/s41586-024-07763-9}{10.1038/s41586-024-07763-9}. Lappalainen et al., Nature, 2024. DOI: \href{https://doi.org/10.1038/s41586-024-07939-3}{10.1038/s41586-024-07939-3}.}
\end{frame}
